% !TEX root=section1.tex
\documentclass[document.tex]{subfiles}

\begin{document}
    \section{A note on notations}
\subsection{Material derivative} \label{material derivative}
The following equations are all equivalent, which can be shown algebraically. 
\begin{align*}
\frac{d}{dt} &= \frac{\partial}{\partial t}+ u \frac{\partial}{\partial x} + v \frac{\partial}{\partial y} + w\frac{\partial}{\partial z}\\
&= \frac{\partial}{\partial t} + \textbf{u}\cdot \nabla \\
&= \frac{\partial}{\partial t} + \textbf{u}\cdot \text{div}() \\
\end{align*}

\subsection{Tensors} \label{tensors}
\begin{align*}
\tau	_{ij}\frac{\partial u_i}{\partial x_j} = \boldsymbol{\tau}:\nabla\textbf{u}
\end{align*}
where 
\begin{align*}
\nabla\mathbf{u} =
\begin{pmatrix}
\dfrac{\partial u}{\partial x} & \dfrac{\partial u}{\partial y} & \dfrac{\partial u}{\partial z} \\[2mm]
\dfrac{\partial v}{\partial x} & \dfrac{\partial v}{\partial y} & \dfrac{\partial v}{\partial z} \\[2mm]
\dfrac{\partial w}{\partial x} & \dfrac{\partial w}{\partial y} & \dfrac{\partial w}{\partial z}
\end{pmatrix}
\end{align*}
The : is called the double contraction operator or the matrix inner product. It contracts two tensors into a scalar. 
Since $\tau$ is the symmetric stress tensor, 

\end{document}